\section{Placement Algorithm}

\vspace{6pt}
Algorithm~\ref{alg:placement} describes the core online scheduling loop.

\begin{tcolorbox}[colback=white,colframe=uwaterlooNavy,title={\textbf{Algorithm 1:} EdgeSched Online Placement}]
\label{alg:placement}
\textbf{Input:} Pending workload \(W = (\mathbf{r}, \ell, \tau, \sigma)\)\\
\textbf{Input:} Node set \(\mathcal{N}\) with current allocations \(\mathcal{A}\)\\
\textbf{Output:} Target node \(n^*\) or reject
\vspace{4pt}
\begin{enumerate}[leftmargin=*,nosep]
    \item Candidate set \(\mathcal{C} \leftarrow \{n \in \mathcal{N} \mid \mathbf{c}_n - \sum_{w \in \mathcal{A}_n} \mathbf{r}_w \geq \mathbf{r}\}\)
    \item \textbf{if} \(\mathcal{C} = \emptyset\) \textbf{then return} reject
    \item For each \(n \in \mathcal{C}\), compute score
          \(s_n = \text{util}(n) + \lambda \cdot \text{colocatePenalty}(n, W)\)
    \item Select \(n^* = \arg\min_{n \in \mathcal{C}} \, s_n\)
    \item \textbf{return} \(n^*\)
\end{enumerate}
\vspace{4pt}
{\small The parameter \(\lambda \in [0, 1]\) trades off utilization against colocation penalty. Higher \(\lambda\) favors dedicated nodes for sensitive workloads.}
\end{tcolorbox}

\vspace{6pt}
The colocation penalty function uses the IAM:

\[
\text{colocatePenalty}(n, W) =
\frac{1}{|\mathcal{A}_n|} \sum_{w \in \mathcal{A}_n} (1 - a_{\tau_w, \tau})
\]

This formulation penalizes nodes that already host workloads with high interference affinity against the incoming workload --- effectively implementing soft anti-affinity without the rigidity of hard constraints.
